\documentclass[11pt]{article}
\usepackage[margin=1in]{geometry}
\usepackage{amsmath}
\usepackage{booktabs}
\usepackage{longtable}
\usepackage{array}
\usepackage{xcolor}
\usepackage{listings}
\usepackage{hyperref}
\hypersetup{colorlinks=true,linkcolor=blue,urlcolor=blue}
\definecolor{codebg}{rgb}{0.96,0.96,0.96}
\lstset{basicstyle=\ttfamily\footnotesize,backgroundcolor=\color{codebg},breaklines=true,frame=single,columns=fullflexible}
\newcommand{\code}[1]{\texttt{#1}}

\title{Pioneer.jl Flanking-Window Features}
\author{Feature reference}
\date{2026-06-24}

\begin{document}
\maketitle

\section{Motivation: wide-isolation-window DIA}
In data-independent acquisition with \emph{wide} quadrupole isolation windows, a single
MS2 window co-isolates many precursors. A confident identification therefore needs evidence
that the precursor's signal is real and not bleed-through from a co-isolated neighbor or
chromatographic interference.

The \textbf{flanking-window features} provide this evidence by contrasting two regions of the
\emph{same} isolation window:
\begin{itemize}
  \item the \textbf{core}: the contiguous run of MS2 cycles, centered on the precursor's
        apex scan, where the precursor passes the search threshold;
  \item the \textbf{flank}: scans of the same isolation window in the cycles immediately
        \emph{flanking} the core (just outside the core's scan range).
\end{itemize}
A genuine precursor has its MS1 and fragment signal concentrated in the core, with fragments
that co-elute; interference tends to smear signal into the flank or to show fragments that do
not track the precursor. The features quantify the core-vs-flank contrast and the
flanking-scan co-elution structure.

\section{Where they are computed (two pipeline stages)}
\begin{enumerate}
  \item \textbf{MainSearch} (\code{SearchMethods/MainSearch/scoring.jl}) --- after
        best-per-precursor selection, \code{\_mainsearch\_flanking\_core\_bounds!} (\code{:811})
        determines the \emph{core} window and the per-precursor loop (\code{:1045--1093})
        writes the intermediate ``core'' columns into the best-PSM table, which are persisted
        in the fold Arrow files:
        \code{flanking\_core\_scan\_min}, \code{flanking\_core\_scan\_max},
        \code{n\_contiguous\_scans}, \code{flanking\_core\_ms1\_m0\_signal},
        \code{flanking\_core\_frag\_signal}.
  \item \textbf{PrecursorScoringSearch} (\code{PrecursorScoringSearch.jl:199} calls
        \code{add\_wide\_window\_features\_to\_fold\_files!}, which calls
        \code{add\_wide\_window\_features\_to\_fold\_file!} once per fold file per MS run, in
        \code{wide\_window\_features.jl}). It reads the core columns, gathers the flanking
        scans, extracts MS1 + fragment intensities from the raw spectra, computes the 11
        flanking feature values in \code{\_wide\_flank\_feature\_values} (\code{:58}), and
        writes them back to the fold Arrow file. The two \code{*\_signal} core columns are
        intermediate and dropped afterward (\code{:893}).
\end{enumerate}
The final feature set (the canonical list \code{FLANKING\_WINDOW\_FEATURES},
\code{model\_config.jl:49}) is then consumed by the second-pass LightGBM in ScoringSearch.

\section{Defining the core (MainSearch)}
\code{\_mainsearch\_flanking\_core\_bounds!} takes the rows of one precursor group (all scans
where the precursor was scored), sorts them by \code{(cycle\_idx, scan\_idx)}, locates the
\emph{best} (selected/apex) scan, and walks outward while each neighbor (a) passes the mask
and (b) is in the immediately adjacent cycle ($\pm 1$):
\begin{lstlisting}
# walk left while passing & cycle is exactly one less
while lo>1 && pass_mask[prev] && cycle[prev]==cycle[cur]-1; lo-=1; end
# walk right while passing & cycle is exactly one greater
while hi<len && pass_mask[next] && cycle[next]==cycle[cur]+1; hi+=1; end
\end{lstlisting}
It returns \code{scan\_min}/\code{scan\_max} (the core scan range), \code{n\_scans}
($=$\code{n\_contiguous\_scans}), and the boundary rows. The core MS1 and fragment
\emph{signals} are then summed over rows whose \code{scan\_idx} lies in
$[\text{scan\_min},\text{scan\_max}]$ (\code{scoring.jl:1086--1093}):
\[
\text{core\_ms1\_m0\_signal} = \!\!\sum_{\text{core scans}}\!\! \text{ms1\_m0}, \qquad
\text{core\_frag\_signal} = \!\!\sum_{\text{core scans}}\!\! \big(\textstyle\sum_{r=1}^{8}\text{frag}_r^{\text{smoothed}}\big).
\]
The 8 core fragment intensities are the smoothed fragment-chromatogram apex values
\code{frag1..8\_smoothed\_intensity} (\code{SMOOTHED\_FRAGMENT\_INTENSITY\_COLUMNS}).

\section{Gathering the flank and extracting signal (PrecursorScoringSearch)}
\begin{itemize}
  \item \code{\_wide\_window\_index\_spectra} (\code{:286}) indexes every MS2 scan by its
        isolation window key \code{(center\_mz, isolation\_width)}, giving the set of scans
        belonging to the same wide window across all cycles.
  \item \code{\_wide\_group\_flank\_window\_scans\_from\_core!} (\code{:420}) collects, for the
        precursor group, the same-window scans in cycles flanking the core (outside
        $[\text{scan\_min},\text{scan\_max}]$, within \code{WIDE\_WINDOW\_CYCLE\_MARGIN}$=3$
        cycles). If there are no flank scans, all features take their zero/default values
        (\code{\_wide\_window\_zero\_feature\_values}).
  \item \code{\_wide\_top\_fragment\_idxs} + \code{\_wide\_group\_fragment\_windows} pick the
        precursor's top-8 library fragments and their m/z tolerance windows.
  \item Scan-centric extraction (\code{\_wide\_fill\_scan\_centric\_ms1\_m0!},
        \code{\_wide\_fill\_scan\_centric\_fragments!}, batched at
        \code{FLANKING\_SCAN\_CENTRIC\_BATCH\_SIZE}$=50000$ groups via
        \code{\_wide\_flush\_scan\_centric\_group\_batch!}) reads each flank scan once and
        fills \code{flank\_ms1\_m0} (length $=$ \#flank scans) and \code{flank\_fragments}
        (\#flank scans $\times$ 8).
\end{itemize}

\section{The feature computations (\code{\_wide\_flank\_feature\_values}, line 58)}
Let the inputs be: core scalars $c_{m0}$ (\code{core\_ms1\_m0\_signal}) and $c_f$
(\code{core\_frag\_signal}); the 8 core apex fragment intensities $a_r$
(\code{core\_fragments}); the flank MS1 vector $\mathbf{m}$ (\code{flank\_ms1\_m0}, one value
per flank scan); and the flank fragment matrix $F$ (\code{flank\_fragments}, scan $\times$
fragment). All intensities are floored at $0$. Define $s_j=\sum_{r} F_{j,r}$ (summed fragment
signal in flank scan $j$), and let $\rho(\cdot,\cdot)$ be Pearson correlation
(\code{\_wide\_window\_pcor}, returns $0$ for $<2$ points or zero variance).

\renewcommand{\arraystretch}{1.35}
\begin{longtable}{>{\raggedright}p{4.4cm} c >{\raggedright\arraybackslash}p{8.2cm}}
\toprule
\textbf{Feature} & \textbf{Type} & \textbf{Definition \& intuition} \\
\midrule
\endhead
\code{flanking\_ms1\_m0\_candidate\_fraction} & Float32 &
$\dfrac{c_{m0}}{c_{m0}+\sum_j \mathbf{m}_j}$ --- fraction of the precursor's MS1 monoisotopic
(m0) signal that sits in the core vs.\ spread into flank scans. $\to 1$ for a clean precursor. \\
\code{flanking\_frag\_candidate\_fraction} & Float32 &
$\dfrac{c_f}{c_f+\sum_j s_j}$ --- same core-vs-flank fraction for summed fragment signal. \\
\code{flanking\_ms1\_frag\_sum\_corr} & Float32 &
$\rho(\mathbf{m},\,\mathbf{s})$ over flank scans --- do the MS1 precursor and the fragment sum
co-vary across the flank? (co-elution sanity check off-apex). \\
\code{flanking\_frag\_corr\_mean} & Float32 &
mean of all pairwise $\rho(F_{:,a},F_{:,b})$ over flank scans, among fragments that have signal
--- average fragment-fragment co-elution in the flank. \\
\code{flanking\_n\_correlated\_fragments} & UInt8 &
count of fragments whose leave-one-out correlation $\rho\!\big(F_{:,r},\,\mathbf{s}-F_{:,r}\big)
> 0.7$ (\code{WIDE\_WINDOW\_CORR\_THRESHOLD}) --- how many fragments track the rest. \\
\code{flanking\_n\_correlated\_fragments\_bitvec\_rank} & UInt16 &
\code{\_bitvec\_pattern\_rank} of the 8-bit mask of \emph{which} fragments exceeded $0.7$ ---
a learned categorical encoding of the correlated-fragment \emph{pattern} (not just the count). \\
\code{flanking\_frag\_corr\_strength} & Float32 &
$\sum_r \max(\min(\rho_r,1),0)$ where $\rho_r$ is the leave-one-out correlation --- continuous
total correlation strength across the 8 fragments. \\
\code{flanking\_frag\_corr\_effective\_n} & Float32 &
$\dfrac{(\sum_r \rho_r^{+})^2}{\sum_r (\rho_r^{+})^2}$ (participation ratio of the positive
leave-one-out corrs) --- effective number of co-eluting fragments. \\
\code{flanking\_frag\_corr\_best\_m0} & Float32 &
$\rho(F_{:,b^\*},\,\mathbf{m})$ where $b^\*$ is the fragment with the highest mean pairwise
correlation (the ``anchor'') --- does the best fragment track the MS1 precursor across the
flank? \\
\code{flanking\_signal\_support} & Float32 &
$\dfrac{\#\{j: \mathbf{m}_j>0 \text{ or } s_j>0\}}{\#\text{flank scans}}$ --- fraction of flank
scans carrying any signal. \\
\code{frag\_apex\_gt2x\_flank\_bitvec\_rank} & UInt16 &
\code{\_bitvec\_pattern\_rank} of the 8-bit mask of fragments whose \emph{core apex}
$a_r \ge 2\times$ their flank average $\bar F_{:,r}$ --- which fragments genuinely peak at the
apex rather than being flat/interference. \\
\midrule
\code{n\_contiguous\_scans} & UInt16 &
(from the core bounds, not \code{\_wide\_flank\_feature\_values}) the number of contiguous
cycles in the core run --- how many cycles the precursor persisted. \\
\bottomrule
\end{longtable}

\noindent The 11 values returned by \code{\_wide\_flank\_feature\_values} plus
\code{n\_contiguous\_scans} make up the 12-entry \code{FLANKING\_WINDOW\_FEATURES} list.

\section{Helper functions}
\begin{longtable}{>{\raggedright}p{6.3cm} >{\raggedright\arraybackslash}p{8.5cm}}
\toprule
\textbf{Function (\code{wide\_window\_features.jl} / \code{scoring.jl})} & \textbf{Role} \\
\midrule
\endhead
\code{\_mainsearch\_flanking\_core\_bounds!} (scoring.jl:811) & Find the contiguous-cycle core
window around the apex; return scan\_min/max, n\_scans, boundary rows. \\
\code{\_wide\_window\_pcor} (:5) & Single-pass Float32 Pearson correlation; $0$ if $<2$ points
or zero variance. \\
\code{\_wide\_candidate\_signal\_fraction} (:32) & $\text{core}/(\text{core}+\text{flank})$,
guarding divide-by-zero. \\
\code{\_wide\_window\_zero\_feature\_values} (:37) & Default feature tuple when there are no
flank scans (fractions still computed from the core signal). \\
\code{\_wide\_flank\_feature\_values} (:58) & Compute the 11 flank features from core +
flank arrays. \\
\code{\_wide\_window\_index\_spectra} (:286) & Index scans by isolation window
\code{(center\_mz, isolation\_width)}. \\
\code{\_wide\_group\_flank\_window\_scans\_from\_core!} (:420) & Collect same-window scans in
cycles flanking the core (margin = 3 cycles). \\
\code{\_wide\_top\_fragment\_idxs} / \code{\_wide\_group\_fragment\_windows} (:453/:511) &
Select the top-8 library fragments and their m/z extraction windows. \\
\code{\_wide\_fill\_scan\_centric\_ms1\_m0!} / \code{\_wide\_fill\_scan\_centric\_fragments!}
(:578/:601) & Read raw-spectra MS1 m0 and fragment intensities at each flank scan (scan-centric,
batched). \\
\code{\_wide\_scatter\_features!} (:706) & Write the computed feature tuple to every row of the
precursor group. \\
\code{\_bitvec\_pattern\_rank} & Map an 8-bit fragment pattern mask to a learned rank
(file-specific table, \code{getBitVecExcessRanks}). \\
\bottomrule
\end{longtable}

\section{Constants}
\code{WIDE\_WINDOW\_CYCLE\_MARGIN}$=3$ (flank reaches up to 3 cycles either side);
\code{WIDE\_WINDOW\_CORR\_THRESHOLD}$=0.7$ (correlated-fragment cutoff);
\code{FLANKING\_SCAN\_CENTRIC\_BATCH\_SIZE}$=50000$ (groups per scan-extraction batch).

\section{Consumption}
The features are appended to the per-fold second-pass PSM Arrow files and become inputs to the
ScoringSearch second-pass LightGBM (via \code{FLANKING\_WINDOW\_FEATURES}). They are most
informative on wide-window instruments (e.g.\ Sciex SWATH), where co-isolation is heaviest.
\end{document}
